\chapter{Experimentation and Evaluation}

Comparing a sparse, linear approach against SotA non-linear models architectures on raw metric scores alone would be an unfair
framing. The evaluation framework therefore operates on three evaluation axes simultaneously:
\begin{itemize}
    \item Axis A — Classification Accuracy: standard discriminative performance metrics (AUROC, F1, etc.)
    \item Axis B — Grounded Localization Quality: volumetric segmentation and localization metrics (DSC, IoU, HIT@K\%)
    \item Axis C — Inherent Explainability: atom-to-voxel traceability, and spatial faithfulness
\end{itemize}

Axis C is the most novel and critical axis of evaluation, and cannot be meaningfully evaluated for the SOTA non-linear
models without expensive post-hoc attribution methods such as Grad-CAM and HiResCAM, that are themselves approximations,
while LC-KSVD provides exact, closed-form localization through its dictionary structure. This asymmetry must be acknowledged,
and the evaluation framework is designed to highlight the unique explainability advantages of LC-KSVD.

\section{Model Registry}
The following registry catalogs every SotA baseline model to be included in the evaluation, their primary task orientation,
and the role they play in the benchmark. Models are grouped by task archetype to clarify the comparison landscape.

\subsection{Classification Baselines}

\subsection{Segmentation/Localization Baselines}

\section{Project Progress}


\section{End Exam} 


%%%%%%%%%%%%%%%%%%%%%%%%%%%%%%%%%%%%%%%%%%%%%%%%%%%%%%%%%%%%%%%%%%%%%%%%%%%%%%%%%%%%%%%%%%%%%%%%%%
%%%%%%%%%%% END OF FILE
%%%%%%%%%%%%%%%%%%%%%%%%%%%%%%%%%%%%%%%%%%%%%%%%%%%%%%%%%%%%%%%%%%%%%%%%%%%%%%%%%%%%%%%%%%%%%%%%%%
